% SUPERSEDED: this section was merged into 01_background.tex for the
% JAMIA submission. It is no longer \input'd from main.tex. Retained
% for version-control history only.
\section{Objectives}
\label{sec:objectives}

We frame our four research questions and the concrete operationalisations
under which we test them.

\textbf{RQ1 (schema operationalisation).}
We compare three nested relation schemas, a 4-label family schema
(S\textsubscript{family}), an 8-label entity-pair schema
(S\textsubscript{pair}), and a 13-label mechanism schema
(S\textsubscript{mech}), and select the operational schema for
downstream evaluation under a pre-committed decision rule:
paired-bootstrap CI on KB-argmax-correctness improvement, with a
Cohen's \(d\) guard against benchmark degradation.

\textbf{RQ2 (training configurations).}
We test three claims: that larger biomedical encoders outperform
their base counterparts; that multi-corpus training outperforms
single-corpus baselines on out-of-domain BC5CDR; and that an
oncology-projected staged schedule outperforms a single-stage
multi-corpus schedule on both external benchmarks.

\textbf{RQ3 (audit-formulation robustness).}
We test whether encoder rankings are stable across three
correctness-aware KB-audit metrics, asking whether the choice of
audit formulation interacts with the choice of encoder.

\textbf{RQ4 (benchmark--KB triangulation).}
We quantify the variance attributable to design levers (schema in
the schema-arm; encoder \(\times\) schedule in the configuration-arm)
in benchmark performance versus KB-surfacing yield, and the rate at
which benchmark-tied configurations preserve their KB ordering. We
treat benchmark and KB performance as two evaluation modalities to
be triangulated rather than as competing rankings.